\documentclass[twocolumn]{aastex631}

\usepackage{amsmath}
\usepackage{multirow}
\usepackage{natbib}
\usepackage{graphicx} 
\usepackage{aas_macros}

\begin{document}

In this paper, we investigated the fundamental question of whether the classical stability of Degenerate Higher-Order Scalar-Tensor (DHOST) theories, which is meticulously engineered to avoid pathological Ostrogradsky ghosts, is robust against quantum corrections. To address this complex problem in a controlled setting, we adopted a strategy guided by symmetry. We first rigorously identified a subclass of Type Ia DHOST theories that possess a field-dependent gauge symmetry and then subjected a representative model from this class to a complete one-loop quantum analysis. Our findings establish a direct and powerful no-go theorem: for this class of theories, the presence of a classical symmetry that becomes anomalous at the quantum level is inextricably linked to a catastrophic loss of quantum stability.

Our methodology involved a three-stage process. First, we derived the integrability conditions required for the existence of the gauge symmetry, thereby carving out a well-defined theoretical laboratory. Second, using the background field method and heat kernel techniques, we computed the one-loop ultraviolet divergences for a simple model within this symmetric class. This calculation revealed that quantum corrections generate local counter-terms quadratic in the curvature, most notably a Weyl-squared invariant. The final and most crucial stage was the analysis of these quantum-generated terms.

The results of this analysis are unambiguous. We first demonstrated that the counter-term Lagrangian does not respect the classical gauge symmetry, giving rise to a quantum anomaly. The effective action, corrected by one-loop effects, is no longer invariant under the transformations that defined the classical theory. More importantly, we showed that this anomaly is not a benign formal feature but a direct signal of a deep physical pathology. The structure of the anomaly-inducing counter-terms, particularly the Weyl-squared term, is fundamentally incompatible with the DHOST framework. These terms cannot be absorbed into a redefinition of the theory's original functions and explicitly violate the algebraic degeneracy conditions that are the cornerstone of the classical theory's stability.

The primary lesson from this work is that the delicate cancellation mechanism responsible for eliminating the Ostrogradsky ghost at the classical level is fragile and can be destroyed by quantum effects. The quantum anomaly serves as a definitive indicator of this breakdown. The reintroduction of the ghost, signaled by the appearance of the Weyl-squared term in the effective action, is not an independent event but a direct consequence of the same quantum effects that break the classical symmetry. This establishes a profound connection between symmetry, anomalies, and the quantum viability of a theory. Consequently, our results elevate the principle of anomaly cancellation from a formal consistency check to a powerful and necessary physical selection principle for constructing consistent theories of modified gravity. While our no-go theorem applies to a specific symmetric subclass, it raises critical questions about the quantum fate of the broader DHOST landscape, suggesting that any classical stability based on fine-tuned cancellations must be protected by a non-anomalous symmetry or another robust quantum mechanism to be considered physically viable.

\end{document}
                